% GENERATED FILE. Edit canonical lessons, then run npm run generate:books.
% canonical-source-hash: fnv1a64:7feedf60a559fda2
% canonical-lessons: HI-C38-aankh, HI-C38-daant, HI-C38-pair, HI-C38-pet

\chapter{Six Words for Where It Hurts}
\label{ch:hi-where-it-hurts}
\hlchaptermodality{38}

\begin{chapteropening}
\textbf{By the end of this chapter:} \emph{I can name six parts of my body, say that one of them is not well, and tell which three are word-for-word cousins of their English translations — and which one has no traceable ancestry at all.}

Say all six body words with the possessive each one demands — \hi{मेरा} \hi{सिर}, \hi{मेरा} \hi{हाथ}, \hi{मेरी} \hi{आँख}, \hi{मेरा} \hi{दाँत}, \hi{मेरा} \hi{पैर}, \hi{मेरा} \hi{पेट} — then say \hi{मेरा} \hi{पेट} \hi{ठीक} \hi{नहीं} \hi{है}, and separate the three exact Indo-European cousins from the two words, \hi{पेट} and \hi{रोटी}, whose histories genuinely stop.
\end{chapteropening}

\section[\hi{आँख}]{\hi{आँख} (āṁkh) — the eye, which is the same word in English}

\label{lesson:HI-C38-aankh}



\begin{quote}
You have a head and a hand. Here is the part of the face — and this one is not merely \emph{related} to its English translation. It is the same word.
\end{quote}

\subsection*{The letters in this word}
\textbf{\hi{आँ}·\hi{ख}} — and the mark on top is an old friend.

\textbf{\hi{आ}} is the standalone long \emph{ā}. Over it sits \textbf{\hi{ँ}}, the \textbf{chandrabindu}, the moon-dot, which sends the vowel up through the nose — exactly what it does in \textbf{\hi{हाँ}}. Then \textbf{\hi{ख}} (\emph{kha}), breathy, with its inherent \emph{a} still attached.

So the word is a nasalised \emph{ā} followed by \emph{kh}: \textbf{āṁkh}, humming at the start.

\begin{grammarlens}[title={Grammar Lens: feminine, and one you will use daily}]
\textbf{\hi{आँख}} (\emph{āṁkh}) = \textquotedblleft{}\textbf{eye},\textquotedblright{} and it is \textbf{feminine} — \textbf{\hi{मेरी} \hi{आँख}}.

Third feminine noun, third consonant ending: \textbf{\hi{चाय}}, \textbf{\hi{किताब}}, \textbf{\hi{आँख}}. If you were still hoping the spelling would rescue you, let this settle it.

And it is worth the effort, because these are the words a doctor asks about. \textbf{\hi{मेरी} \hi{आँख} \hi{ठीक} \hi{नहीं} \hi{है}} — \textquotedblleft{}my eye is not well.\textquotedblright{}
\end{grammarlens}

\begin{cousinweb}
\textbf{\hi{आँख}} comes from Sanskrit \textbf{\hi{अक्षि}} (\emph{akṣi}), \textquotedblleft{}eye\textquotedblright{} — and \emph{akṣi} comes from the root \emph{\textbf{h\textsubscript{3}ek\textsuperscript{w}-}}, which is one of the securest words in the whole family.

Its descendants:

\begin{itemize}
\raggedright
  \item Latin \textbf{oculus} — inside English \textbf{ocular}, \textbf{monocle}, \textbf{binoculars}.
  \item Old English \textbf{ēage}, which is English \textbf{eye} itself.
  \item Lithuanian \textbf{akìs}, Armenian \textbf{akn}, and Greek \textbf{ósse}, the old dual form meaning \textquotedblleft{}the two eyes.\textquotedblright{}
\end{itemize}

A caution on the Greek, because the obvious candidate is a trap: \textbf{ophthalmós} is \emph{usually} connected here, but its first part has never been explained, and serious opinion treats it as a pre-Greek word. Cite \textbf{ósse} and stay honest.

\textbf{\hi{हाथ}} had to travel a long, battered road from \emph{hasta}. \textbf{\hi{आँख}} barely travelled at all — \emph{akṣi} to \emph{āṁkh} is a short, clean walk, and at the far end of it an English speaker is saying a version of the same word.
\end{cousinweb}

\subsection*{Your turn}
Say these aloud:

\begin{itemize}
\raggedright
  \item \emph{āṁkh}, then \emph{merī āṁkh}
  \item the three body words you now have — \emph{sir}, \emph{hāth}, \emph{āṁkh}
  \item \emph{hāṁ} and \emph{āṁkh} — the same moon-dot, the same hum
  \item \emph{akṣi}, \emph{oculus}, \emph{eye} — one root, three languages
\end{itemize}

\begin{tcolorbox}[breakable,colback=teal!4,colframe=teal!35!black,title={Before you move on}]
Is \textbf{\hi{आँख}} masculine or feminine? (\textbf{Feminine} — \emph{merī āṁkh}.) Which Greek word is the safe cousin, and which is the trap? (\textbf{ósse} is safe; \textbf{ophthalmós} is unexplained.) What does the \textbf{\hi{ँ}} over the \textbf{\hi{आ}} do, and where did you first meet it? (\textbf{Nasalises} the vowel — on \textbf{\hi{हाँ}}.)
\end{tcolorbox}
\section[danta]{\hi{दाँत} (dāṁt) — the tooth, and the gender trap of this chapter}

\label{lesson:HI-C38-daant}



\begin{quote}
Built like the last word, spelled like the last word, and the opposite gender. This is the pair learners get wrong most often.
\end{quote}

\subsection*{The letters in this word}
\textbf{\hi{दाँ}·\hi{त}} — three known parts and no new letter at all.

\textbf{\hi{द}} (\emph{da}) takes \textbf{\hi{ा}}, and the \textbf{\hi{ँ}} rides above the head-line for the nasal hum — the same \textbf{chandrabindu} that turns \textbf{\hi{ह}} into \textbf{\hi{हाँ}}, \textquotedblleft{}yes.\textquotedblright{} Then \textbf{\hi{त}} (\emph{ta}).

Both \textbf{\hi{द}} and \textbf{\hi{त}} are made at the \textbf{teeth} — which is the whole point of Devanagari's mouth-map. The word for \emph{tooth} is spelled with the two consonants your tongue makes \emph{against your teeth}. Say \emph{dāṁt} and feel where it happens.

\begin{grammarlens}[title={Grammar Lens: the trap}]
\textbf{\hi{दाँत}} (\emph{dāṁt}) = \textquotedblleft{}\textbf{tooth},\textquotedblright{} and it is \textbf{masculine} — \textbf{\hi{मेरा} \hi{दाँत}}.

Now put the two side by side. \textbf{\hi{आँख}} and \textbf{\hi{दाँत}}: same long \emph{ā}, same chandrabindu, same single final consonant, both parts of the body, both everyday. \textbf{\hi{आँख}} is feminine. \textbf{\hi{दाँत}} is masculine.

There is no rule underneath that. There is no reasoning that gets you from one to the other. This is the pair to hold up whenever you are tempted to guess: \textbf{\hi{मेरी} \hi{आँख}}, \textbf{\hi{मेरा} \hi{दाँत}}.
\end{grammarlens}

\begin{cousinweb}
\textbf{\hi{दाँत}} is Sanskrit \textbf{\hi{दन्त}} (\emph{danta}), from \emph{\textbf{h\textsubscript{3}dónts}} — and like \emph{akṣi} this is one of the family's best-attested words.

\begin{itemize}
\raggedright
  \item Latin \textbf{dēns}, stem \emph{dentis} — English \textbf{dentist}, \textbf{dental}, \textbf{trident} (a \textquotedblleft{}three-tooth\textquotedblright{}).
  \item Greek \textbf{odoús} — English \textbf{odontology}.
  \item Old English \textbf{tōþ}, which is \textbf{tooth}.
\end{itemize}

Two body parts in a row, both plain cousins of their English translations. That is not a coincidence about eyes and teeth; it is what a core inherited vocabulary looks like from the inside.

And the older reading of \emph{\textbf{h\textsubscript{3}dónts}} makes it a participle of \emph{\textbf{h\textsubscript{3}ed-}}, \textquotedblleft{}to bite\textquotedblright{} — so a tooth is \textquotedblleft{}\textbf{the biting one},\textquotedblright{} which is what \textbf{\hi{खाना}} needs it for. Same move as \textbf{\hi{दूध}}, \textquotedblleft{}the milked one\textquotedblright{}: name the thing after what happens.
\end{cousinweb}

\subsection*{Your turn}
Say these aloud:

\begin{itemize}
\raggedright
  \item \emph{dāṁt}, then \emph{merā dāṁt}
  \item the trap pair, twice — \emph{merī āṁkh, merā dāṁt}; again
  \item \emph{hāṁ}, then \emph{dāṁt} — the same moon-dot doing the same job
  \item \emph{danta}, \emph{dēns}, \emph{odoús}, \emph{tooth} — one word, four languages
  \item the two \textquotedblleft{}named for the action\textquotedblright{} words — \emph{dugdha}, the milked; the tooth, the biting one
\end{itemize}

\begin{tcolorbox}[breakable,colback=teal!4,colframe=teal!35!black,title={Before you move on}]
\textbf{\hi{आँख}} and \textbf{\hi{दाँत}} look alike — do they share a gender? (\textbf{No} — \emph{merī āṁkh}, \emph{merā dāṁt}.) Which English words carry the Latin form? (\textbf{Dentist}, \textbf{dental}, \textbf{trident}.) What does the older reading make a tooth? (\textbf{The biting one} — from a root meaning \textquotedblleft{}to bite.\textquotedblright{})
\end{tcolorbox}
\section[\hi{पैर}]{\hi{पैर} (pair) — the foot, and the ancestor it does not have}

\label{lesson:HI-C38-pair}



\begin{quote}
Down from the head, past the hand, to what carries you. And this word's obvious ancestor turns out to belong to a different Hindi word.
\end{quote}

\begin{grammarlens}[title={Grammar Lens: masculine, and the sentence that uses it}]
\textbf{\hi{पैर}} (\emph{pair}) = \textquotedblleft{}\textbf{foot},\textquotedblright{} and it is \textbf{masculine} — \textbf{\hi{मेरा} \hi{पैर}}.

The \textbf{\hi{ै}} is a single mātrā for one vowel, \emph{ai} — not \emph{a} then \emph{i}. One sign, one sound, sitting above the head-line on \textbf{\hi{प}}.

Now the sentence you would actually say to a doctor, using the \textbf{\hi{नहीं}} you already have: \textbf{\hi{मेरा} \hi{पैर} \hi{ठीक} \hi{नहीं} \hi{है}} — \textquotedblleft{}\textbf{my foot is not well}.\textquotedblright{} Watch the \textbf{\hi{नहीं}} sit directly before \textbf{\hi{है}}, which is where Hindi puts its negation.
\end{grammarlens}

\begin{cousinweb}
Here is the obvious guess, and it is wrong. Sanskrit has \textbf{\hi{पाद}} (\emph{pāda}), \textquotedblleft{}foot\textquotedblright{} — long \emph{ā}, famous, the source of English's borrowed \emph{pada-} words. It looks like it must be the parent of \emph{pair}. It is not.

\textbf{\hi{पाँव}} (\emph{pāṁv}), Hindi's \emph{other} everyday word for foot, is the descendant of \textbf{\hi{पाद}}. \textbf{\hi{पैर}} comes from the shorter Sanskrit stem \textbf{\hi{पद्}} (\emph{pad-}), through Prakrit \emph{paya}, picking up a Middle Indo-Aryan \textbf{-ra-} ending on the way.

So Hindi inherited the same body part \textbf{twice}, down two different stems of the same word — and both survive, side by side, meaning the same thing.

Further back both are \emph{\textbf{ped-}}, \textquotedblleft{}to walk, to tread\textquotedblright{}:

\begin{itemize}
\raggedright
  \item Latin \textbf{pēs}, stem \emph{pedis} — \textbf{pedal}, \textbf{pedestrian}, \textbf{pedicure}.
  \item Greek \textbf{poús}, stem \emph{podós} — \textbf{podium}, and the eight of an \textbf{octopus}.
  \item Old English \textbf{fōt}, which is \textbf{foot}.
\end{itemize}

Three body words in three lessons, and every one of them is the same word your own language uses. \textbf{\hi{हाथ}} had to be argued for through a sound change. These three simply \emph{are} the cousins.
\end{cousinweb}

\subsection*{Your turn}
Say these aloud:

\begin{itemize}
\raggedright
  \item \emph{pair}, then \emph{merā pair}
  \item the four body words with their possessives — \emph{merā hāth, merī āṁkh, merā dāṁt, merā pair}
  \item \emph{merā pair ṭhīk nahīṁ hai} — and hear where the \emph{nahīṁ} sits
  \item which Hindi word came from \emph{pāda}, and which from \emph{pad-} (\emph{pāṁv}; \emph{pair})
\end{itemize}

\begin{tcolorbox}[breakable,colback=teal!4,colframe=teal!35!black,title={Before you move on}]
Does \textbf{\hi{पैर}} come from \textbf{\hi{पाद}}? (\textbf{No} — from \textbf{\hi{पद्}}; \textbf{\hi{पाँव}} is the \emph{pāda} one.) Which English words carry the Latin and Greek forms? (\textbf{Pedal}, \textbf{pedestrian}; \textbf{podium}, \textbf{octopus}.) Say \textquotedblleft{}my foot is not well,\textquotedblright{} and say where \textbf{\hi{नहीं}} goes. (\emph{Merā pair ṭhīk nahīṁ hai} — immediately before \textbf{\hi{है}}.)
\end{tcolorbox}
\section[\hi{पेट}]{\hi{पेट} (peṭ) — the belly, and a trail that genuinely stops}

\label{lesson:HI-C38-pet}



\begin{quote}
Three clean cousins in a row. This one breaks the run — and the break is worth more than another neat answer.
\end{quote}

\subsection*{The letters in this word}
\textbf{\hi{पे}·\hi{ट}} — one known letter, one known mātrā, one consonant to feel.

\textbf{\hi{प}} takes \textbf{\hi{े}}, the mātrā that perches on the head-line and turns the built-in \emph{a} into \emph{e} — the second of the two vowel signs you practised on \textbf{\hi{नाम}} and \textbf{\hi{मेरा}}.

Then \textbf{\hi{ट}}: a \textbf{retroflex} \emph{ṭ}, tongue-tip curled back, not the \textbf{\hi{त}} of \textbf{\hi{दाँत}} made flat against the teeth. Say \emph{dāṁt}, then \emph{peṭ}.

\begin{grammarlens}[title={Grammar Lens: masculine, and the whole set together}]
\textbf{\hi{पेट}} (\emph{peṭ}) = \textquotedblleft{}\textbf{stomach, belly},\textquotedblright{} and it is \textbf{masculine} — \textbf{\hi{मेरा} \hi{पेट}}.

Now you can answer a doctor, and exactly one of the six takes \textbf{\hi{मेरी}}: \emph{merā sir}, \emph{merā hāth}, \textbf{merī āṁkh}, \emph{merā dāṁt}, \emph{merā pair}, \emph{merā peṭ} — feminine company for \textbf{\hi{कुर्सी}} and \textbf{\hi{किताब}}.

The sentence, with \textbf{\hi{नहीं}} in its usual place just before \textbf{\hi{है}}: \textbf{\hi{मेरा} \hi{पेट} \hi{ठीक} \hi{नहीं} \hi{है}} — which, said in India, means what it says and rather more.
\end{grammarlens}

\begin{cousinweb}
\textbf{\hi{पेट}} goes back to Prakrit \textbf{\hi{पेट्ट}} (\emph{peṭṭa}) — and there the clean line ends.

The mainstream account treats it as a \textbf{Dravidian} word taken into Indo-Aryan; compare reconstructed Dravidian \textbf{*poṭṭ-}. Turner's comparative dictionary does not name Dravidian, but neither does he derive \emph{peṭṭa} from any Sanskrit word: he only remarks that it sits close to Indo-Aryan words for \textquotedblleft{}\textbf{basket}.\textquotedblright{} That resemblance is a note, not a derivation, and it is often quoted as one.

So: \textbf{the origin of \hi{पेट} is unsettled}, with Dravidian leading — the honesty \textbf{\hi{रोटी}} needed, one lesson away from \textbf{\hi{पानी}}, whose story is exact. A language is made of what people said, and some of that came from neighbours speaking something else.

There is even a fingerprint: the \textbf{retroflex} consonants, the \textbf{\hi{ट}} you just curled your tongue for, are the feature of Indian Indo-Aryan most often credited to long contact with Dravidian.
\end{cousinweb}

\subsection*{Your turn}
Say these aloud:

\begin{itemize}
\raggedright
  \item all six — \emph{merā sir, merā hāth, merī āṁkh, merā dāṁt, merā pair, merā peṭ}
  \item \emph{merā peṭ ṭhīk nahīṁ hai}; then \emph{merī āṁkh ṭhīk nahīṁ hai}
  \item every feminine noun you own — \emph{kursī}, \emph{chāy}, \emph{kitāb}, \emph{āṁkh}
  \item the three exact ancestries — \emph{āṁkh}, \emph{dāṁt}, \emph{pair}; then the two unsettled — \emph{peṭ}, \emph{roṭī}
\end{itemize}

\begin{tcolorbox}[breakable,colback=teal!4,colframe=teal!35!black,title={Before you move on}]
Where does \textbf{\hi{पेट}} come from? (\textbf{Unsettled} — Prakrit \emph{peṭṭa}, most likely a \textbf{Dravidian} loan; the \textquotedblleft{}basket\textquotedblright{} link is a resemblance.) Of your six body words, which takes \textbf{\hi{मेरी}}? (\textbf{\hi{आँख}}.) Which earlier word had the same honest dead end? (\textbf{\hi{रोटी}}, beside \textbf{\hi{पानी}}.)
\end{tcolorbox}
